\section{2-dimensional coherence diagrams}
\label{app:coh2}

The diagrams below expand the seven hexagonal and four square obligations in Figures~\ref{fig:arch-coh-basic}--\ref{fig:arch-coh-refl}. Vertices are composites of boundary operations, edges are instances of frame equalities, and the central double arrow denotes equality of the two directed paths with common endpoints. Each such 2-dimensional coherence follows from $\UIP$ for frames.

Here $\delta$ abbreviates restriction, and $\varepsilon^\downarrow$ and $\varepsilon^\uparrow$ abbreviate Below and Above degeneracy. Edge labels use the same abbreviations in coherence names: for example, $\delta\text{-}\varepsilon^\downarrow\text{-}\mathrm{inf}$ denotes $\relRestrReflBelowInf$.

The displayed composites use the relative indices of the construction, with remaining-layer index $k$; dimension superscripts, frame and painting arguments, and dependent transports are implicit. An edge is whiskered by the surrounding unchanged operations. Composition is read from right to left, and paths are followed in the direction of the arrows. The labels ``zero'' and ``successor'' refer to the layer cases of Section~\ref{sec:coherences}. Within a successor diagram, $q$ and $r$ are the indices of the recursive premises, after removing that successor. The face labels $\alpha,\beta,\gamma$ range over $\nu$.

\begingroup
% Path diagrams adapted from the Bonak higher-coherence companion.
\newcommand{\del}[2]{\delta^{#2}_{#1}}
\newcommand{\epsdmark}{\varepsilon^\downarrow}
\newcommand{\epsumark}{\varepsilon^\uparrow}
\newcommand{\epsd}[1]{\varepsilon^\downarrow_{#1}}
\newcommand{\epsu}[1]{\varepsilon^\uparrow_{#1}}
\newcommand{\comp}{\mathbin{\circ}}

\newcommand{\HC}[2]{#1\!\left(#2\right)}
\newcommand{\HCdd}[4]{\HC{\delta\text{-}\delta}{#1,#2,#3,#4}}
\newcommand{\HCdeid}[2]{\HC{\delta\text{-}\epsdmark\text{-}\mathrm{id}}{#1,#2}}
\newcommand{\HCedinf}[3]{\HC{\delta\text{-}\epsdmark\text{-}\mathrm{inf}}{#1,#2,#3}}
\newcommand{\HCedsup}[3]{\HC{\delta\text{-}\epsdmark\text{-}\mathrm{sup}}{#1,#2,#3}}
\newcommand{\HCeusup}[3]{\HC{\delta\text{-}\epsumark\text{-}\mathrm{sup}}{#1,#2,#3}}
\newcommand{\HCeded}[2]{\HC{\epsdmark\text{-}\epsdmark}{#1,#2}}
\newcommand{\HCeued}[2]{\HC{\epsdmark\text{-}\epsumark}{#1,#2}}
\newcommand{\HCeueu}[2]{\HC{\epsumark\text{-}\epsumark}{#1,#2}}
\newcommand{\termone}[1]{#1}
\newcommand{\termthree}[3]{#1 \comp #2 \comp #3}
\tikzset{
  coh edge/.style={->,>=Latex,thick},
  coh term/.style={align=center,font=\normalsize,inner sep=1.5pt},
  coh label/.style={fill=white,inner sep=1pt,font=\normalsize}
}

\newcommand{\HexTopEdges}[3]{%
  \draw[coh edge] (L) -- node[coh label,above left] {$#1$} (UL);
  \draw[coh edge] (UL) -- node[coh label,above] {$#2$} (UR);
  \draw[coh edge] (UR) -- node[coh label,above right] {$#3$} (R);
}

\newcommand{\HexBottomEdges}[3]{%
  \draw[coh edge] (L) -- node[coh label,below left] {$#1$} (LL);
  \draw[coh edge] (LL) -- node[coh label,below] {$#2$} (LR);
  \draw[coh edge] (LR) -- node[coh label,below right] {$#3$} (R);
}

\newcommand{\HexTermCohDiagram}[8]{%
  \begin{tikzpicture}[baseline=(current bounding box.center)]
    \node[coh term] (L) at (0,0) {$#1$};
    \node[coh term] (UL) at (4.5,3) {$#2$};
    \node[coh term] (UR) at (11,3) {$#3$};
    \node[coh term] (R) at (15.5,0) {$#4$};
    \node[coh term] (LR) at (11,-3) {$#5$};
    \node[coh term] (LL) at (4.5,-3) {$#6$};
    \HexTopEdges#7
    \HexBottomEdges#8
    \node at (7.75,0) {$\Downarrow$};
  \end{tikzpicture}%
}

\newcommand{\SquareTermCohDiagram}[8]{%
  \begin{tikzpicture}[baseline=(current bounding box.center)]
    \node[coh term] (TL) at (0,3.3) {$#1$};
    \node[coh term] (TR) at (10.5,3.3) {$#2$};
    \node[coh term] (BL) at (0,0) {$#3$};
    \node[coh term] (BR) at (10.5,0) {$#4$};
    \draw[coh edge] (TL) -- node[coh label,above] {$#5$} (TR);
    \draw[coh edge] (TR) -- node[coh label,right] {$#6$} (BR);
    \draw[coh edge] (TL) -- node[coh label,left] {$#7$} (BL);
    \draw[coh edge] (BL) -- node[coh label,below] {$#8$} (BR);
    \node at (5.25,1.65) {$\Downarrow$};
  \end{tikzpicture}%
}

\newcommand{\SquareTermCohDiagramRevB}[8]{%
  \begin{tikzpicture}[baseline=(current bounding box.center)]
    \node[coh term] (TL) at (0,3.3) {$#1$};
    \node[coh term] (TR) at (10.5,3.3) {$#2$};
    \node[coh term] (BL) at (0,0) {$#3$};
    \node[coh term] (BR) at (10.5,0) {$#4$};
    \draw[coh edge] (TL) -- node[coh label,above] {$#5$} (TR);
    \draw[coh edge] (TR) -- node[coh label,right] {$#6$} (BR);
    \draw[coh edge] (TL) -- node[coh label,left] {$#7$} (BL);
    \draw[coh edge] (BR) -- node[coh label,below] {$#8$} (BL);
    \node at (5.25,1.65) {$\Downarrow$};
  \end{tikzpicture}%
}


Equation~\eqref{eq:coh1-recap} recalls the eight frame-coherence families in the notation used on the diagram edges. These are the construction-oriented equalities of Tables~\ref{tab:cohframe-types} and~\ref{tab:cohframe-dgn}, with arguments and transports suppressed. In particular, cancellation is oriented from the identity to restriction after degeneracy. Each coherence name is followed by its equality and index bounds; the diagrams compare composites of these proofs. All direction indices are nonnegative and face labels belong to $\nu$.

\begin{equation}
  \label{eq:coh1-recap}
  \begin{array}{l@{\;}c@{\;}l@{\;}c@{\;}l@{\qquad}l}
    \HCdd{r}{q}{\gamma}{\beta} & : & \del{q}{\beta}\comp\del{r}{\gamma} & = & \del{r}{\gamma}\comp\del{q+1}{\beta} & r\leq q\leq k     \\[0.4em]
    \HCdeid{r}{\beta}          & : & \id                                & = & \del{r}{\beta}\comp\epsd{r}          & r\leq k           \\[0.4em]
    \HCedinf{r}{q}{\beta}      & : & \epsd{q}\comp\del{r}{\beta}        & = & \del{r}{\beta}\comp\epsd{q+1}        & r\leq q\leq k     \\[0.4em]
    \HCedsup{r}{q}{\beta}      & : & \epsd{q}\comp\del{r}{\beta}        & = & \del{r+1}{\beta}\comp\epsd{q}        & q\leq r\leq k     \\[0.4em]
    \HCeusup{r}{q}{\beta}      & : & \epsu{q}\comp\del{r}{\beta}        & = & \del{r}{\beta}\comp\epsu{q}          & q\leq p,\;r\leq k \\[0.4em]
    \HCeded{r}{q}              & : & \epsd{q}\comp\epsd{r}              & = & \epsd{r+1}\comp\epsd{q}              & q\leq r\leq k     \\[0.4em]
    \HCeued{r}{q}              & : & \epsu{q}\comp\epsd{r}              & = & \epsd{r}\comp\epsu{q}                & q\leq p,\;r\leq k \\[0.4em]
    \HCeueu{r}{q}              & : & \epsu{r+1}\comp\epsu{q}            & = & \epsu{q}\comp\epsu{r}                & q\leq r\leq p
  \end{array}
\end{equation}

\subsection{\texorpdfstring{$\relRestrRestr$}{restr-restr}}

\begin{figure}[H]
  \centering
  \resizebox{\linewidth}{!}{\HexTermCohDiagram
    {\termthree{\del{q}{\beta}}{\del{r}{\gamma}}{\del{s}{\alpha}}}
    {\termthree{\del{q}{\beta}}{\del{s}{\alpha}}{\del{r+1}{\gamma}}}
    {\termthree{\del{s}{\alpha}}{\del{q+1}{\beta}}{\del{r+1}{\gamma}}}
    {\termthree{\del{s}{\alpha}}{\del{r+1}{\gamma}}{\del{q+2}{\beta}}}
    {\termthree{\del{r}{\gamma}}{\del{s}{\alpha}}{\del{q+2}{\beta}}}
    {\termthree{\del{r}{\gamma}}{\del{q+1}{\beta}}{\del{s}{\alpha}}}
    {{\HCdd{s}{r}{\alpha}{\gamma}}{\HCdd{s}{q}{\alpha}{\beta}}{\HCdd{r+1}{q+1}{\gamma}{\beta}}}
    {{\HCdd{r}{q}{\gamma}{\beta}}{\HCdd{s}{q+1}{\alpha}{\beta}}{\HCdd{s}{r}{\alpha}{\gamma}}}
}\par\medskip
  $s\leq r\leq q\leq k$\par
  \caption{2-dimensional coherence for $\relRestrRestr$. The hexagon compares the two orders of commuting three restrictions.}
  \label{fig:coh2-dd}
\end{figure}

\subsection{\texorpdfstring{$\relRestrReflId$}{restr-refl-id}}

\begin{figure}[H]
  \centering
  \resizebox{\linewidth}{!}{\SquareTermCohDiagramRevB
    {\termone{\del{s}{\alpha}}}
    {\termthree{\del{r}{\beta}}{\epsd{r}}{\del{s}{\alpha}}}
    {\termthree{\del{s}{\alpha}}{\del{r+1}{\beta}}{\epsd{r+1}}}
    {\termthree{\del{r}{\beta}}{\del{s}{\alpha}}{\epsd{r+1}}}
    {\HCdeid{r}{\beta}}
    {\HCedinf{s}{r}{\alpha}}
    {\HCdeid{r+1}{\beta}}
    {\HCdd{s}{r}{\alpha}{\beta}}
}\par\medskip
  $s\leq r\leq k$\par
  \caption{2-dimensional coherence for $\relRestrReflBelowId$. The square relates cancellation to the mixed and restriction coherences.}
  \label{fig:coh2-id}
\end{figure}

\subsection{\texorpdfstring{$\relRestrReflInf$}{restr-refl-inf}}

\begin{figure}[H]
  \centering
  \resizebox{\linewidth}{!}{\HexTermCohDiagram
    {\termthree{\epsd{q}}{\del{r}{\beta}}{\del{s}{\alpha}}}
    {\termthree{\epsd{q}}{\del{s}{\alpha}}{\del{r+1}{\beta}}}
    {\termthree{\del{s}{\alpha}}{\epsd{q+1}}{\del{r+1}{\beta}}}
    {\termthree{\del{s}{\alpha}}{\del{r+1}{\beta}}{\epsd{q+2}}}
    {\termthree{\del{r}{\beta}}{\del{s}{\alpha}}{\epsd{q+2}}}
    {\termthree{\del{r}{\beta}}{\epsd{q+1}}{\del{s}{\alpha}}}
    {{\HCdd{s}{r}{\alpha}{\beta}}{\HCedinf{s}{q}{\alpha}}{\HCedinf{r+1}{q+1}{\beta}}}
    {{\HCedinf{r}{q}{\beta}}{\HCedinf{s}{q+1}{\alpha}}{\HCdd{s}{r}{\alpha}{\beta}}}
}\par\medskip
  $s\leq r\leq q\leq k$\par
  \caption{2-dimensional coherence for $\relRestrReflBelowInf$. The hexagon moves a Below degeneracy past two restrictions.}
  \label{fig:coh2-inf}
\end{figure}

\subsection{\texorpdfstring{$\relRestrReflSup$}{restr-refl-sup}}

\begin{figure}[H]
  \centering
  \resizebox{\linewidth}{!}{\HexTermCohDiagram
    {\termthree{\epsd{q}}{\del{r}{\beta}}{\del{s}{\alpha}}}
    {\termthree{\epsd{q}}{\del{s}{\alpha}}{\del{r+1}{\beta}}}
    {\termthree{\del{s}{\alpha}}{\epsd{q+1}}{\del{r+1}{\beta}}}
    {\termthree{\del{s}{\alpha}}{\del{r+2}{\beta}}{\epsd{q+1}}}
    {\termthree{\del{r+1}{\beta}}{\del{s}{\alpha}}{\epsd{q+1}}}
    {\termthree{\del{r+1}{\beta}}{\epsd{q}}{\del{s}{\alpha}}}
    {{\HCdd{s}{r}{\alpha}{\beta}}{\HCedinf{s}{q}{\alpha}}{\HCedsup{r+1}{q+1}{\beta}}}
    {{\HCedsup{r}{q}{\beta}}{\HCedinf{s}{q}{\alpha}}{\HCdd{s}{r+1}{\alpha}{\beta}}}
}\par\medskip
  $s\leq q\leq r\leq k$\par
  \caption{2-dimensional coherence for $\relRestrReflBelowSup$. The hexagon compares the two orders of applying the restriction and mixed coherences.}
  \label{fig:coh2-sup-below}
\end{figure}

\begin{figure}[H]
  \centering
  \resizebox{\linewidth}{!}{\HexTermCohDiagram
    {\termthree{\epsu{q}}{\del{r}{\beta}}{\del{s}{\alpha}}}
    {\termthree{\epsu{q}}{\del{s}{\alpha}}{\del{r+1}{\beta}}}
    {\termthree{\del{s}{\alpha}}{\epsu{q}}{\del{r+1}{\beta}}}
    {\termthree{\del{s}{\alpha}}{\del{r+1}{\beta}}{\epsu{q}}}
    {\termthree{\del{r}{\beta}}{\del{s}{\alpha}}{\epsu{q}}}
    {\termthree{\del{r}{\beta}}{\epsu{q}}{\del{s}{\alpha}}}
    {{\HCdd{s}{r}{\alpha}{\beta}}{\HCeusup{s}{q}{\alpha}}{\HCeusup{r+1}{q}{\beta}}}
    {{\HCeusup{r}{q}{\beta}}{\HCeusup{s}{q}{\alpha}}{\HCdd{s}{r}{\alpha}{\beta}}}
}\par\medskip
  $q\leq p,\quad s\leq r\leq k$\par
  \caption{2-dimensional coherence for $\relRestrReflAboveSup\text{ (successor)}$. The successor Above stage uses a hexagon.}
  \label{fig:coh2-sup-above}
\end{figure}

\begin{figure}[H]
  \centering
  \resizebox{\linewidth}{!}{\SquareTermCohDiagram
    {\termone{\del{r}{\beta}}}
    {\termthree{\del{s}{\alpha}}{\epsd{s}}{\del{r}{\beta}}}
    {\termthree{\del{r}{\beta}}{\del{s}{\alpha}}{\epsd{s}}}
    {\termthree{\del{s}{\alpha}}{\del{r+1}{\beta}}{\epsd{s}}}
    {\HCdeid{s}{\alpha}}
    {\HCedsup{r}{s}{\beta}}
    {\HCdeid{s}{\alpha}}
    {\HCdd{s}{r}{\alpha}{\beta}}
}\par\medskip
  $s\leq r\leq k$\par
  \caption{2-dimensional coherence for $\relRestrReflAboveSup\text{ (zero)}$. At the transition to Above, the cancellation coherence closes a square.}
  \label{fig:coh2-sup-zero}
\end{figure}

\subsection{\texorpdfstring{$\relReflRefl$}{refl-refl}}

\begin{figure}[H]
  \centering
  \resizebox{\linewidth}{!}{\HexTermCohDiagram
    {\termthree{\epsd{q}}{\epsd{r}}{\del{s}{\alpha}}}
    {\termthree{\epsd{q}}{\del{s}{\alpha}}{\epsd{r+1}}}
    {\termthree{\del{s}{\alpha}}{\epsd{q+1}}{\epsd{r+1}}}
    {\termthree{\del{s}{\alpha}}{\epsd{r+2}}{\epsd{q+1}}}
    {\termthree{\epsd{r+1}}{\del{s}{\alpha}}{\epsd{q+1}}}
    {\termthree{\epsd{r+1}}{\epsd{q}}{\del{s}{\alpha}}}
    {{\HCedinf{s}{r}{\alpha}}{\HCedinf{s}{q}{\alpha}}{\HCeded{r+1}{q+1}}}
    {{\HCeded{r}{q}}{\HCedinf{s}{q}{\alpha}}{\HCedinf{s}{r+1}{\alpha}}}
}\par\medskip
  $s\leq q\leq r\leq k$\par
  \caption{2-dimensional coherence for $\relReflBelowBelow$. The hexagon compares commuting the degeneracies before or after moving the restriction through them.}
  \label{fig:coh2-refl-bb}
\end{figure}

\begin{figure}[H]
  \centering
  \resizebox{\linewidth}{!}{\HexTermCohDiagram
    {\termthree{\epsu{q}}{\epsd{r}}{\del{s}{\alpha}}}
    {\termthree{\epsu{q}}{\del{s}{\alpha}}{\epsd{r+1}}}
    {\termthree{\del{s}{\alpha}}{\epsu{q}}{\epsd{r+1}}}
    {\termthree{\del{s}{\alpha}}{\epsd{r+1}}{\epsu{q}}}
    {\termthree{\epsd{r}}{\del{s}{\alpha}}{\epsu{q}}}
    {\termthree{\epsd{r}}{\epsu{q}}{\del{s}{\alpha}}}
    {{\HCedinf{s}{r}{\alpha}}{\HCeusup{s}{q}{\alpha}}{\HCeued{r+1}{q}}}
    {{\HCeued{r}{q}}{\HCeusup{s}{q}{\alpha}}{\HCedinf{s}{r}{\alpha}}}
}\par\medskip
  $q\leq p,\quad s\leq r\leq k$\par
  \caption{2-dimensional coherence for $\relReflBelowAbove\text{ (successor)}$. The mixed successor stage uses a hexagon.}
  \label{fig:coh2-refl-ba}
\end{figure}

\begin{figure}[H]
  \centering
  \resizebox{\linewidth}{!}{\SquareTermCohDiagram
    {\termone{\epsd{r}}}
    {\termthree{\del{s}{\alpha}}{\epsd{s}}{\epsd{r}}}
    {\termthree{\epsd{r}}{\del{s}{\alpha}}{\epsd{s}}}
    {\termthree{\del{s}{\alpha}}{\epsd{r+1}}{\epsd{s}}}
    {\HCdeid{s}{\alpha}}
    {\HCeded{r}{s}}
    {\HCdeid{s}{\alpha}}
    {\HCedinf{s}{r}{\alpha}}
}\par\medskip
  $s\leq r\leq k$\par
  \caption{2-dimensional coherence for $\relReflBelowAbove\text{ (zero)}$. The transition from Below--Below uses a square.}
  \label{fig:coh2-refl-ba-zero}
\end{figure}

\begin{figure}[H]
  \centering
  \resizebox{\linewidth}{!}{\HexTermCohDiagram
    {\termthree{\epsu{r+1}}{\epsu{q}}{\del{s}{\alpha}}}
    {\termthree{\epsu{r+1}}{\del{s}{\alpha}}{\epsu{q}}}
    {\termthree{\del{s}{\alpha}}{\epsu{r+1}}{\epsu{q}}}
    {\termthree{\del{s}{\alpha}}{\epsu{q}}{\epsu{r}}}
    {\termthree{\epsu{q}}{\del{s}{\alpha}}{\epsu{r}}}
    {\termthree{\epsu{q}}{\epsu{r}}{\del{s}{\alpha}}}
    {{\HCeusup{s}{q}{\alpha}}{\HCeusup{s}{r+1}{\alpha}}{\HCeueu{r}{q}}}
    {{\HCeueu{r}{q}}{\HCeusup{s}{r}{\alpha}}{\HCeusup{s}{q}{\alpha}}}
}\par\medskip
  $q\leq r\leq p,\quad s\leq k$\par
  \caption{2-dimensional coherence for $\relReflAboveAbove\text{ (successor)}$. The Above--Above successor stage uses a hexagon.}
  \label{fig:coh2-refl-aa}
\end{figure}

\begin{figure}[H]
  \centering
  \resizebox{\linewidth}{!}{\SquareTermCohDiagramRevB
    {\termone{\epsu{r}}}
    {\termthree{\epsu{r}}{\del{s}{\alpha}}{\epsd{s}}}
    {\termthree{\del{s}{\alpha}}{\epsd{s}}{\epsu{r}}}
    {\termthree{\del{s}{\alpha}}{\epsu{r}}{\epsd{s}}}
    {\HCdeid{s}{\alpha}}
    {\HCeusup{s}{r}{\alpha}}
    {\HCdeid{s}{\alpha}}
    {\HCeued{s}{r}}
}\par\medskip
  $r\leq p,\quad s\leq k$\par
  \caption{2-dimensional coherence for $\relReflAboveAbove\text{ (zero)}$. The transition from the mixed stage uses a square.}
  \label{fig:coh2-refl-aa-zero}
\end{figure}
\endgroup
